\documentclass[11pt]{article}
\usepackage{amsmath}
\title{One frame correction, many pairs of acquisition and use timestamps}
\date{26 November 2026}
\begin{document}
\maketitle
\paragraph{Transfer from the selected source studies.}
The preceding projection study established that a shared gauge
correction between the scan and target timestamps need not cancel.
The selected monitor source also distinguishes the repeated use of
one measurement from its acquisition. This study gives an exact replay
of the resulting observation interface under a declared scalar
correction model. Timestamp provenance, the bound on one pending
correction, and the declared exact unperturbed projections are
inputs. This is not an identification of a real correction history.

\paragraph{Joint model and algorithm.}
There is one correction at an unknown instant $\tau$ in a closed
interval $C$, with post-correction signed residual $z$ in a closed
interval $J$. The correction is already effective at equality.
For use $i$ its acquisition is $a_i$, its evaluation cut is $u_i$,
with $a_i\le u_i$, and its unperturbed projected returns are $x_{ij}$.
Use $i$ adds the shift $z$ to its point list if
$a_i<\tau\le u_i$, and adds zero otherwise. This is the one-step
translation case of the two-time reference transformation,
with the sign of the correction absorbed in $z$.
Successive uses may cite the same acquisition. Their deterministic
projected lists can differ. They advance the native evaluation
counter and a separate distinct-acquisition corroboration check.
The latter does not certify capture order or safety by itself.
The modeled identity and known event clocks must actually be
provided to apply either guarantee.

Split the correction-time interval at each acquisition and use stamp
and the jump interval at the polygon incidence boundaries for the
current point list. Retain singleton boundaries as well as open
cells. The use vote is constant on each rectangular product cell:
the time coordinate specifies whether the pair crosses, and the
shift coordinate specifies the membership count if it does. Starting
from the known pair of state machines, a new row splits the partition,
inherits the previous product state on each subcell, and applies
the same current hit or miss to both machines. Every concrete
$(\tau,z)$ follows exactly its own product cell. Induction then
proves exactness over traces without predicting later uses.
With $T$ unique time boundaries and $K$ unique shift boundaries
there are at most $(2T-1)(2K-1)$ cells. This is a state-space count,
not timing or storage guarantees on a deployed controller.
The direct comparison reconstructs the complete partition for each
prefix and evaluates concrete histories from scratch. Its results
match 1{,}162 prefixes and 40{,}572 cell representatives.

\paragraph{Sharp distinction at touching windows.}
Use the overlap-sensitive half-open stamp intervals $(0,2]$ and
$(2,4]$, threshold one, and a polygon accepting $0\le x<1$ at
$y=1/2$. A common correction-shift in $[7/8,1]$ moves a return at
$-3/4$ into the polygon. Each use could hit, but no single correction
can place both uses inside. With an initially inactive latch,
a two-consecutive-hit trigger is therefore impossible, even though an
independent per-window model accepts it. Changing only the second
acquisition to time 1 makes a shared trigger possible precisely
because $(1,2]$ is nonempty. Conversely, two in-polygon returns
at $1/4$ move outside if their uses cross that same correction.
Two separate misses, and hence release from an active two-miss
latch with two corroborating acquisitions, remain impossible for the
touching windows but become possible when they overlap. An exact
timestamp equality is consequential. These are exact synthetic
event-order contrasts, not assurance about a real stop or release.

Unbounded drift, multiple coupled corrections, uncertain stamp order,
occlusion and nonshared errors need different models. In particular,
clipping a current interval after forgetting the old association
between a correction and its effects does not recreate that history.
A factored baseline retaining the declared joint history can obtain
the same result. No new family or architecture acceptance is claimed.
\end{document}
